\textbf{Risk Model for ILP algo}

There are 2 sets of securities :

• outright set (size \emph{M}) e.g. 7 near/far futures \emph{M=14}

• tradable set (size \emph{N≥M}) which include outright set e.g. 7
near/far futures 7 spread \emph{M=21}

\textbf{Part 1. Functor FCovar}

The risk model of ILP also starts with finding the covariance matrix of
daily return among 14 outrights, which include 7 EGB near futures plus 7
EGB far futures, hence \emph{M=14}, for example :

• FGBSU6 and FGBSZ6

• FGBMU6 and FGBMZ6

• FGBLU6 and FGBLZ6

• FGBXU6 and FGBXZ6

• FBTSU6 and FBTSZ6

• FBTPU6 and FBTPZ6

• FOATU6 and FOATZ6

There is a daily python job in \emph{qfilib} to calculate the 14 by 14
covariance matrix from market data on each day retrospectively, which
then performs PCA and extracts the 3 principal eigen vectors,
corresponding eigen values and 14 residuals on diagonal of covariance.
The \emph{qfilib} job then saves the PCA result in KDB, which will be
read by another unity process called \emph{egb\_refdata}, and published
to stew folder. ILP algo then subscribes the stew folder for the PCA
result from stew folder. Here is what it looks like in HPD. For
simplicity, it shows 7 near outrights only, yet in production there are
14 outrights. The yellow box matrix is not about PCA, it is about
liquidity adjusted VaR, please refer to next doc.

\includegraphics[width=6.26806in,height=1.85in]{media/image1.png}

Dimension reduction with rank3-PCA is a noise-removal process by
retaining major yield curve risk :

• PC1 = yield curve level shift

• PC2 = yield curve tilt

• PC3 = yield curve curvature

Residual ensures the covariance values on the diagonal are lossless
after dimension reduction. ILP algo recontructs covariance matrix by
functor \hl{FCovar} using equation :

\(Cov\) \emph{=} \(Q\Lambda Q^{T} + R\)

where \(Q\) \emph{= 14 × 3 matrix for eigen vector \textbf{(green box)}}

\(\Lambda\) \emph{= 3 × 3 diagonal matrix for eigen value \textbf{(blue
box)}}

\(R\) \emph{= 14 × 14 diagonal matrix for residuals \textbf{(red box)}}

ILP algo does not consider tradable set in \hl{FCovar}. Offline
retrospective covariance analysis is done with outright set only. Thus
algos with different tradable sets, but same outright set can share the
same \hl{FCovar} without recalibrating its own \hl{FCovar}. Numbers in
yellow box are for risk adjusted VaR (see next doc).

\textbf{Part 2. Functor FRisk}

Risk model of ILP algo takes the following input :

• covariance of outright set \emph{(rank3-PCA approximation)}

• position of outright set \emph{(all tradables are decomposed into
outrights)}

then calculates the following risk items :

• portfolio risk (variance)

• implied position of tradable set

• gradient of portfolio risk in outright space called bias / shading

• delta position for outright-wise risk minimization called maxRiskRed

\emph{\textbf{2a. Position in tradable domain / outright domain}}

Let adopt the following notation :

\(X_{trd}\) \emph{position of tradable set, size N}

\(X_{out}\) \emph{position of tradable set, decomposed into outrights,
size M}

\(P_{trd} = \{ p_{0},p_{1},\ldots,p_{N - 1}\}\) \emph{price of tradable
set, size N}

\(P_{out} = \{ p_{0}',p_{1}',\ldots,p_{M - 1}'\}\) \emph{price of
outright set, size M}

\emph{where}
\(\{ p_{0}',\ldots,p_{M - 1}'\} \subseteq \{ p_{0},\ldots,p_{N - 1}\}\)

We consider \(X_{out}\) instead of \(X_{trd}\) for risk model, because
covariance matrix is calibrated overnight using outrights only, avoid
including strategies like spreads or butterflies which can be different
for different algos. First of all, portfolio of tradable set should have
the same value as portfolio of outright set, if the outright set
perfectly replicates tradable set :

\(X_{trd}^{T}P_{trd}\) = \(X_{out}^{T}P_{out}\)

Given replication matrix \emph{S} which is \emph{N} by \emph{M}
(\emph{S} is transpose of that in code), by no arbitrage pricing :

\(P_{trd}\) = \(SP_{out}\) \emph{no arbitrage pricing}

\(S_{n}\) = \(\lbrack\ldots, + 1, - 1,\ldots\rbrack\) \emph{suppose nth
row of S is a spread}

\(S_{n}\) = \(\lbrack\ldots, - 1, + 2, - 1,\ldots\rbrack\) \emph{suppose
nth row of S is a butterfly}

\(S_{n}\) = \(\lbrack\ldots, + 1, + 1, + 1, + 1\ldots\rbrack\)
\emph{suppose nth row of S is a pack}

from which, we can derive the relation between \(X_{trd}\) and
\(X_{out}\) :

\(X_{trd}^{T}P_{trd}\) = \(X_{out}^{T}P_{out}\)

\emph{→} \(X_{trd}^{T}SP_{out}\) = \(X_{out}^{T}P_{out}\) \emph{applied
no arbitrage pricing}

\emph{→} \(X_{trd}^{T}S\) = \(X_{out}^{T}\)

\emph{→} \(X_{out}\) = \(S^{T}X_{trd}\)

To summarise (please note that \emph{S} matrix in code may be transpose
of convention in here) :

\(X_{trd}^{T}P_{trd}\) = \(X_{out}^{T}P_{out}\) \emph{(eq1) portfolio
replication}

\(P_{trd}\) = \(SP_{out}\) \emph{(eq2) no arbitrage pricing}

\begin{quote}
\(X_{out}\) = \(S^{T}X_{trd}\) \emph{(eq3) decompose tradable set into
outright set}
\end{quote}

\emph{\textbf{\hfill\break
}}

\emph{\textbf{2b. Covariance matrix and portfolio risk}}

\(C_{trd}\) = \(Cov\lbrack\mathrm{\Delta}P_{trd}\rbrack\)

= \(Cov\left\lbrack S\mathrm{\Delta}P_{out} \right\rbrack\) \emph{since}
\({\mathrm{\Delta}P}_{trd} = \mathrm{\Delta}\left( SP_{out} \right) = S\mathrm{\Delta}P_{out}\)

= \(SCov\left\lbrack \mathrm{\Delta}P_{out} \right\rbrack S^{T}\)

= \(SC_{out}S^{T}\) \emph{(eq4) \hl{implemented in FRisk::updateCovar}}

\(risk\) = \(X_{trd}^{T}C_{trd}X_{trd}\)

= \(X_{trd}^{T}SC_{out}S^{T}X_{trd}\) \emph{apply eq4}

= \({(S^{T}X_{trd})}^{T}C_{out}S^{T}X_{trd}\)

= \(X_{out}^{T}C_{out}X_{out}\) \emph{apply eq3 \hl{implemented in
FRisk::updateRisk}}

Portfolio risk can be expressed in tradable domain or in outrights
domain, both are equivalent. \emph{Eq6} is used instead of \emph{eq5},
because only \(C_{out}\) (not \(C_{trd}\)) is available in overnight
calibration process. That makes calibration easier as it involves
outright set only. \emph{In short,} \emph{no} \(C_{trd}\) \emph{in KDB}.

\(risk\) = \(X_{trd}^{T}C_{trd}X_{trd}\) \emph{(eq5)}

= \(X_{out}^{T}C_{out}X_{out}\) \emph{(eq6)}

\emph{\textbf{2c. Objective function (i.e. functor}} FRisk
\emph{\textbf{value) and derivative}}

The objective function is portfolio risk plus liquidity cost of all
tradable securities. Liquidity cost can be different for different
securities, proportional to position size regardless of sign of
position, i.e. the same liquidity cost for long or short position of the
same amount. It is transaction cost plus spread the desk needs to pay to
offset a position.

\(f(X_{trd})\) =
\(risk + 2\sum_{n}^{}{{lq}_{n}\left| x_{trd,n} \right|}\)

\begin{quote}
=
\(X_{trd}^{T}C_{trd}X_{trd} + 2\sum_{n}^{}{{lq}_{n}\left| x_{trd,n} \right|}\)
\end{quote}

= \(X_{trd}^{T}C_{trd}X_{trd} + 2{lq}^{T}\left| X_{trd} \right|\)
\emph{where lq is column vector with size N}

= \(X_{out}^{T}C_{out}X_{out} + 2{lq}^{T}\left| X_{trd} \right|\)
\emph{where \textbar\textbar{} is elementwise absolute}

Functor FRisk also calculates gradient of \(J(X_{trd})\), we take
derivatives with respect to \(X_{trd}\), instead of \(X_{out}\), because
we are going to hedge with tradable sets, not with trading outrights.

\(\frac{1}{2}\ \frac{df(X_{trd})}{dX_{trd}}\) =
\(\frac{1}{2}\frac{d}{dX_{trd}}(X_{out}^{T}C_{out}X_{out}) + \frac{d}{dX_{trd}}{lq}^{T}\left| X_{trd} \right|\)

=
\(\frac{1}{2}\frac{d}{dX_{trd}}(X_{trd}^{T}SC_{out}S^{T}X_{trd}) + \frac{d}{dX_{trd}}{lq}^{T}\left| X_{trd} \right|\)
\emph{apply eq3}

= \(SC_{out}S^{T}X_{trd} + \sum_{n}^{}{{lq}_{n}\mu_{n}}\)

= \(SC_{out}X_{out} + \sum_{n}^{}{{lq}_{n}\mu_{n}}\) \emph{apply eq3
again}

where \(\mu_{n}\) = \(\left\lbrack \begin{matrix}
 + 1 & if & x_{n} > 0 \\
 - 1 & if & x_{n} < 0 \\
0 & if & x_{n} = 0
\end{matrix} \right.\ \)

In ffwk code, we define :

FRisk::value() = \(f(X_{trd})\) =
\(X_{out}^{T}C_{out}X_{out} + 2{lq}^{T}\left| X_{trd} \right|\)

FRisk::m\_bias = \(\frac{1}{2}\ \frac{df(X_{trd})}{dX_{trd}}\) =
\(SC_{out}X_{out} + \sum_{n}^{}{{lq}_{n}\mu_{n}}\)

\emph{\textbf{\hfill\break
}}

\emph{\textbf{2d. Marginal risk reduction (per security risk
reduction)}}

If we hedge by \({\mathrm{\Delta}X}_{trd}\), objective function becomes
\(J(X_{trd} + \mathrm{\Delta}X_{trd})\), the minimum is trivial, perfect
hedge can give minimum with zero objective function (also zero
liquidation cost) :

\(\min_{\mathrm{\Delta}X_{trd}}{f(X_{trd} + \mathrm{\Delta}X_{trd})}\) =
\(f(0)\) \emph{when} \(\mathrm{\Delta}X_{trd} = - X_{trd}\)

= \(0\)

Yet, a perfect hedge for all tradables simultaneously is impossible in
practice. We cannot synchronize the hedges for all securities, as it
assumes infinity market liquidity and zero trading latency. In practice,
we have to do marginal variance reduction iteratively by trading
security one at each time. Consider the case with 2 correlated
securities having position \(X_{trd} = \lbrack - 2, - 1\rbrack\), the
ideal hedge will be buy 2 and buy 1 shares respectively, which is not
achievable, hence we need to do it iteratively. In each step, we do not
completely offset the position of a security, because a smaller
portfolio varaince can be achieved with non zero position due to the
correlation between the tradables.

\includegraphics[width=3.33402in,height=2.65953in]{media/image2.png}

Now consider minimizing objective \(f(X_{trd})\) by trading only with
tradable \emph{k}, that is, by changing :

\(x_{trd,m}\) \emph{→} \(x_{trd,m}\) \emph{for m ≠ k}

\(x_{trd,k}\) → \(x_{trd,k} + \mathrm{\Delta}x\) \emph{for m = k (use
symbol} \(\mathrm{\Delta}x\) \emph{instead of}
\(\mathrm{\Delta}x_{trd,k}\) \emph{for simplicity)}

\emph{hence} \(X_{trd}\) → \(X_{trd} + \mathrm{\Delta}x\delta_{k}\)

\emph{where} \(\delta_{k}\) \emph{= 1 at k-th position, 0 otherwise
(column vector of size N)}

Recall objective \(f(X_{trd})\) as follows, and we denote
\(f(X_{trd} + \mathrm{\Delta}x)\) as \emph{1D} quadratic
\(f(\mathrm{\Delta}x)\) for simplicity :

\begin{quote}
\(f(X_{trd})\) =
\(X_{out}^{T}C_{out}X_{out} + 2{lq}^{T}\left| X_{trd} \right|\)
\end{quote}

= \(X_{trd}^{T}SC_{out}S^{T}X_{trd} + 2{lq}^{T}\left| X_{trd} \right|\)

\(f(\mathrm{\Delta}x)\) =
\({(X_{trd} + \mathrm{\Delta}x\delta_{k})}^{T}SC_{out}S^{T}(X_{trd} + \mathrm{\Delta}x\delta_{k}) + 2\sum_{m \neq k}^{}{{lq}_{m}\left| x_{trd,m} \right|} + 2{lq}_{k}\left| x_{trd,k} + \mathrm{\Delta}x \right|\)

=
\((\mathrm{\Delta}x)^{2}\left( \delta_{k}^{T}SC_{out}S^{T}\delta_{k} \right) + 2\mathrm{\Delta}x\left( \delta_{k}^{T}SC_{out}S^{T}X_{trd} \right) +\)

\(\underset{const}{\overset{X_{trd}^{T}SC_{out}S^{T}X_{trd} + 2\sum_{m \neq k}^{}{{lq}_{m}\left| x_{trd,m} \right|}}{︸}} + 2{lq}_{k}\left| x_{trd,k} + \mathrm{\Delta}x \right|\)

=
\((\mathrm{\Delta}x)^{2}\left( \delta_{k}^{T}C_{trd}\delta_{k} \right) + 2\mathrm{\Delta}x\left( \delta_{k}^{T}SC_{out}X_{out} \right) + 2{lq}_{k}\left| x_{trd,k} + \mathrm{\Delta}x \right| + const\)

=
\((\mathrm{\Delta}x)^{2}C_{trd,k,k} + 2\mathrm{\Delta}x\left( S_{k}C_{out}X_{out} \right) + 2{lq}_{k}\left| x_{trd,k} + \mathrm{\Delta}x \right| + const\)

\(f'(\mathrm{\Delta}x)\) =
\(2\mathrm{\Delta}xC_{trd,k,k} + 2\left( S_{k}C_{out}X_{out} \right) + 2{\mu_{k}lq}_{k}\)

\(f'\left( {\mathrm{\Delta}x}_{opt} \right)\) =
\(2{\mathrm{\Delta}x}_{opt}C_{trd,k,k} + 2\left( S_{k}C_{out}X_{out} \right) + 2{\mu_{k}lq}_{k}\)
= 0

\({\mathrm{\Delta}x}_{opt}\) \emph{=}
\(- \ \frac{S_{k}C_{out}X_{out} + {\mu_{k}lq}_{k}}{C_{trd,k,k}}\)

\emph{where} \(S_{k}\) = \(\delta_{k}^{T}S\) = \emph{k-th row of S (row
vector of size M)}

\(C_{trd,k,k}\) = \(\delta_{k}^{T}C_{trd}\delta_{k}\) = \emph{(k,k)
element of tradable covariance}

\(\mu_{k}\) = \(\left\lbrack \begin{matrix}
 + 1 & if & x_{trd,k} + {\mathrm{\Delta}x}_{opt} > 0 \\
 - 1 & if & x_{trd,k} + {\mathrm{\Delta}x}_{opt} < 0 \\
\lbrack - 1, + 1\rbrack & if & x_{trd,k} + {\mathrm{\Delta}x}_{opt} = 0
\end{matrix} \right.\ \) = \emph{slope of}
\(\left| x_{trd,k} + {\mathrm{\Delta}x}_{opt} \right|\)

Applying LASSO \emph{\textbf{soft threshold operator}} :

\({\mathrm{\Delta}x}_{opt}\) = \(\left\lbrack \begin{matrix}
 - \ \frac{S_{k}C_{out}X_{out} + {lq}_{k}}{C_{trd,k,k}} & if & {\mathrm{\Delta}x}_{opt} > - x_{trd,k} \\
 - \ \frac{S_{k}C_{out}X_{out} - {lq}_{k}}{C_{trd,k,k}} & if & {\mathrm{\Delta}x}_{opt} < - x_{trd,k} \\
 - x_{trd,k} & if & otherwise
\end{matrix} \right.\ \)

which means :

\({\mathrm{\Delta}x}_{opt}\) = \(\left\lbrack \begin{matrix}
{\mathrm{\Delta}x}_{opt}^{+} & if & - x_{trd,k} < {\mathrm{\Delta}x}_{opt}^{+} < {\mathrm{\Delta}x}_{opt}^{-} & where & {\mathrm{\Delta}x}_{opt}^{+} = - \ \frac{S_{k}C_{out}X_{out} + {lq}_{k}}{C_{trd,k,k}} \\
{\mathrm{\Delta}x}_{opt}^{-} & if & {{\mathrm{\Delta}x}_{opt}^{+} < \mathrm{\Delta}x}_{opt}^{-} < - x_{trd,k} & where & {\mathrm{\Delta}x}_{opt}^{-} = - \ \frac{S_{k}C_{out}X_{out} - {lq}_{k}}{C_{trd,k,k}} \\
 - x_{trd,k} & if & {\mathrm{\Delta}x}_{opt}^{+} < - x_{trd,k} < {\mathrm{\Delta}x}_{opt}^{-} & & 
\end{matrix} \right.\ \)

since liquidation cost \({lq}_{k}\) is always positive :

\({\mathrm{\Delta}x}_{opt}^{+}\ \) \textless{}
\({\mathrm{\Delta}x}_{opt}^{-}\)

In other words, if 100\% hedge of tradable \emph{k} (i.e.
\(- x_{trd,k}\)) lies within the band
\({\lbrack\mathrm{\Delta}x}_{opt}^{+},{\mathrm{\Delta}x}_{opt}^{-}\rbrack\),
then optimal is 100\% hedge, otherwise if \(- x_{trd,k}\) lies outside
the band, we just clamp it to the nearest boundary. The 100\% hedge
means the benefit of \emph{\textbf{marginal risk reduction}} \emph{(i.e.
per security risk reducetion)} is less than the liquidation cost of that
security, then we would rather 100\% offset it.

In ffwk code, we define :

FRisk::m\_targetVar{[}n{]} = \(C_{trd,n,n}\)
\(\forall n \in \lbrack 1,N\rbrack\)

FRisk::m\_maxRiskRed{[}n{]} = \({\mathrm{\Delta}x}_{opt}\) \emph{for
tradable n} \(\forall n \in \lbrack 1,N\rbrack\)

\emph{\textbf{2e. Implied tradable position (from outright position)}}

In part 2a, we know how to decompose tradable position into outright
position, yet there is no way to imply \(X_{trd} \in \mathcal{R}^{N}\)
from \(X_{out} \in \mathcal{R}^{M}\), where \(N \geq M\) (since
tradables contains outrights). There are multiple possible tradable
positions, that can be decomposed into the same outright position.
Therefore \hl{a proper design is} that the functor FRisk \hl{takes
tradable position as input}, and convert it to outright position inside
the functor, it will then have accurate position for both domains.
However for some reasons, FRisk takes outright position as input
instead, hence it needs to imply tradable position.

\emph{\ul{Approach 1 -- Imply tradable position using least square}}

Assuming that :

1. tradables do not contain outright and

2. number of tradables is less than number of outrights (i.e.
\(M \geq N\)), then

\begin{quote}
\(S^{T}X_{trd}\) \textasciitilde{} \(X_{out}\) \emph{LS model}
\end{quote}

\(\widehat{X_{trd}}\) =
\(\underset{X}{argmin}{{(S^{T}X - X_{out})}^{T}(S^{T}X - X_{out})}\)
\emph{LS error}

\textbf{=} \({({SS}^{T})}^{- 1}(SX_{out})\) \emph{unfortunately} FRisk
\emph{is not doing this}

\emph{\ul{Approach 2 -- Projection of outright position on n-th row of
S}}

The current implementation inside FRisk is :

\(X_{trd,n}\) = \(\frac{S_{n}X_{out}}{\left| S_{n} \right|^{2}}\) =
\emph{projection of} \(X_{out}\) \emph{on} \(S_{n}\)

\emph{where} \(S_{n}\) = \(\delta_{n}^{T}S\) = \emph{n-th row of S (row
vector of size M)}

The physical meaning of projection is that, if the entire book (i.e. all
outright positions) are held as \emph{n-th} tradable only, how many
shares of tradable \emph{m} would there be. However that does not make
much sense. Consider 1 unit spread \emph{M6-U6}, 1 unit spread
\emph{U6-Z6}, plus 0 unit spread \emph{M6-Z6}, hence
\(X_{out} = {\lbrack 1,0, - 1\rbrack}^{T}\), implied position for the 3
spread contracts are :

\(X_{trd,0}\) =
\(\frac{\lbrack 1, - 1,0\rbrack X_{out}}{|1, - 1,0|^{2}}\) =
\(\frac{1}{2}\)

\(X_{trd,1}\) =
\(\frac{\lbrack 0,1, - 1\rbrack X_{out}}{|0,1, - 1|^{2}}\) =
\(\frac{1}{2}\)

\(X_{trd,2}\) =
\(\frac{\lbrack 1,0, - 1\rbrack X_{out}}{|1,0, - 1|^{2}}\) = \(1\)

In ffwk code, we define :

FRisk::m\_impliedPos{[}n{]} =
\(\frac{S_{n}X_{out}}{\left| S_{n} \right|^{2}}\)
\(\forall n \in \lbrack 1,N\rbrack\)

\emph{\textbf{2f. Benchmark}}

Functor FRisk also calculates a few benchmark items. Added them in the
following summary.

FRisk::value() = \(f(X_{trd})\) =
\(X_{out}^{T}C_{out}X_{out} + 2{lq}^{T}\left| X_{trd} \right|\)

FRisk::m\_bias = \(\frac{1}{2}\ \frac{df(X_{trd})}{dX_{trd}}\) =
\(SC_{out}X_{out} + \sum_{n}^{}{{lq}_{n}\mu_{n}}\)

FRisk::m\_maxRiskRed{[}n{]} = \(\left\lbrack \begin{matrix}
 - \ \frac{S_{n}C_{out}X_{out} + {lq}_{n}}{C_{trd,n,n}} & if & {\mathrm{\Delta}x}_{opt} > - x_{trd,k} \\
 - \ \frac{S_{n}C_{out}X_{out} - {lq}_{n}}{C_{trd,n,n}} & if & {\mathrm{\Delta}x}_{opt} < - x_{trd,k} \\
 - x_{trd,n} & if & otherwise
\end{matrix} \right.\ \) \(\forall n \in \lbrack 1,N\rbrack\)

FRisk::m\_impliedPos{[}n{]} =
\(\frac{S_{n}X_{out}}{\left| S_{n} \right|^{2}}\)
\(\forall n \in \lbrack 1,N\rbrack\)

FRisk::m\_targetVar{[}n{]} = \(C_{trd,n,n}\)
\(\forall n \in \lbrack 1,N\rbrack\)

FRisk::m\_targetVol{[}n{]} = \(\sqrt{C_{trd,n,n}}\)
\(\forall n \in \lbrack 1,N\rbrack\)

FRisk::m\_benchRisk = \({A_{bench}^{2}C}_{trd,bench,bench}\)

FRisk::m\_benchQty{[}n{]} =
\(A_{bench}\sqrt{\frac{C_{trd,bench,bench}}{C_{trd,n,n}}}\)
\(\forall n \in \lbrack 1,N\rbrack\)

FRisk::m\_targetRange{[}n{]} =
\(\sqrt{\frac{R_{bench}{A_{bench}^{2}C}_{trd,bench,bench}}{C_{trd,n,n}}}\)
\(\forall n \in \lbrack 1,N\rbrack\)

\textbf{\hfill\break
}

\textbf{Appendix : Least Absolute Shrinking and Selection Operator
(LASSO)}

LASSO is a common techique when we add L1 penalty to least square
regression problem. Lets review.

For a typical least square problem, L1 and L2 penalty are (\emph{A is
N×M, B is N×1 and X is M×1}) :

\(\min_{X}{((AX - B)^{T}(AX - B) + \lambda\sum_{m}^{}\left| x_{m} \right|)}\)
\emph{least square with L1 penalty, a.k.a. LASSO}

\(\min_{X}{((AX - B)^{T}(AX - B) + \lambda X^{T}X)}\) \emph{least square
with L2 penalty, a.k.a. ridge regression}

Both above objective functions have 2 parts, the first part for
minimising sum of error square, the second part for forcing each
parameter close to zero. The contours for \emph{L2} penalty are
concentric circles (spheres) with origin as the centre, while the
contours for \emph{L1} penalty are diamonds in 2D, octahedrons in 3D
with minimum at the centre.

\includegraphics[width=4.9375in,height=2.206in]{media/image3.png}

\emph{L2} penalty is differentiable, hence its solution is obvious :

\(f(X)\) = \((AX - B)^{T}(AX - B) + \lambda X^{T}X\)

\(\nabla f(X)\) = \({2A}^{T}(AX - B) + 2\lambda X\)

\(\nabla f\left( X_{opt} \right)\) =
\({2(A}^{T}AX_{opt} - A^{T}B) + 2\lambda X_{opt}\) = \(0\)

\(X_{opt}\) = \({{(A}^{T}A + \lambda I)}^{- 1}A^{T}B\)

\emph{L1} penalty is not differentiable, need to consider different
cases. For each dimension, there are 3 cases :

\(\left\lbrack \begin{matrix}
x_{m} > 0 \\
x_{m} < 0 \\
x_{m} = 0
\end{matrix} \right.\ \) \(\forall m \in \lbrack 1,M\rbrack\)

Therefore there will be \(3^{M}\) cases. Consider one out of the
\(3^{M}\) cases, with
\(\mu = {{\lbrack\mu}_{0},\mu_{1},\ldots,\mu_{M - 1}\rbrack}^{T}\) so
that

\(\left\lbrack \begin{matrix}
\mu_{m} = + 1 & if & x_{m} > 0 \\
\mu_{m} = - 1 & if & x_{m} < 0 \\
\mu_{m} = 0 & if & x_{m} = 0
\end{matrix} \right.\ \)

\(f(X)\) =
\((AX - B)^{T}(AX - B) + \lambda\sum_{m}^{}\left| x_{m} \right|\)

= \((AX - B)^{T}(AX - B) + \lambda\mu^{T}X\)

\(\nabla f(X)\) = \({2A}^{T}(AX - B) + \lambda\mu\)

\(\nabla f\left( X_{opt} \right)\) =
\({2(A}^{T}AX_{opt} - A^{T}B) + \lambda\mu\) = \(0\)

\(X_{opt}\) = \({{(A}^{T}A)}^{- 1}(A^{T}B - \frac{1}{2}\lambda\mu)\)

After solving \(X_{opt}\), we need to verify if it is consistent with
assumption in \(\mu\), if yes, take this as a solution, if not, abandone
it, go for next combination \(\mu\), however this approach needs to
retry \(3^{M}\) cases, which is infeasible. Instead it can be
implemented iteratively by fixing all except one dimensions of parameter
unchanged, optimize by solve for that dimension \(x_{k}\), then move on
to next, repeat until convergence.

Iteration is done as following. Given current estimate \(X\), we update
dimension \emph{k} from \(x_{k}\) to
\(x_{k} + {\mathrm{\Delta}x}_{k}\delta_{k}\), while keeping all other
dimensions unchanged. Find the value of scalar
\({\mathrm{\Delta}x}_{k}\) that minimises \emph{L1} cost.

\(x_{m}\) \emph{→} \(x_{m}\) \emph{for m ≠ k}

\(x_{k}\) \emph{→} \(x_{k} + \mathrm{\Delta}x_{k}\) \emph{for m = k}

\emph{hence} \(X\) \emph{→} \(X + {\mathrm{\Delta}x}_{k}\delta_{k}\)

\emph{where} \(\delta_{k}\) \emph{= 1 at k-th position, 0 otherwise
(column vector of size M)}

Objective function
\(f\left( X + {\mathrm{\Delta}x}_{k}\delta_{k} \right)\) becomes a
\emph{1D} quadratic function \(f({\mathrm{\Delta}x}_{k})\) :

\(f\left( {\mathrm{\Delta}x}_{k} \right)\) =
\(\left( A(X + {\mathrm{\Delta}x}_{k}\delta_{k}) - B \right)^{T}\left( A(X + {\mathrm{\Delta}x}_{k}\delta_{k}) - B \right) + \lambda\sum_{m \neq k}^{}\left| x_{m} \right| + \lambda\left| x_{k} + {\mathrm{\Delta}x}_{k} \right|\)

=
\(\left( {\mathrm{\Delta}x}_{k}{A\delta}_{k} + (AX - B) \right)^{T}\left( {\mathrm{\Delta}x}_{k}{A\delta}_{k} + (AX - B) \right) + \lambda\sum_{m \neq k}^{}\left| x_{m} \right| + \lambda\left| x_{k} + {\mathrm{\Delta}x}_{k} \right|\)

=
\(\left( {\mathrm{\Delta}x}_{k} \right)^{2}\left( {A\delta}_{k} \right)^{T}\left( {A\delta}_{k} \right) + {2\mathrm{\Delta}x}_{k}\left( {A\delta}_{k} \right)^{T}(AX - B) +\)

\(\underset{const}{\overset{(AX + B)^{T}(AX - B) + \lambda\sum_{m \neq k}^{}\left| x_{m} \right|}{︸}} + \lambda\left| x_{k} + {\mathrm{\Delta}x}_{k} \right|\)

=
\(\left( {\mathrm{\Delta}x}_{k} \right)^{2}\left( {A\delta}_{k} \right)^{T}\left( {A\delta}_{k} \right) + {2\mathrm{\Delta}x}_{k}\left( {A\delta}_{k} \right)^{T}(AX - B) + \lambda\left| x_{k} + {\mathrm{\Delta}x}_{k} \right| + const\)

=
\(\left( {\mathrm{\Delta}x}_{k} \right)^{2}(\delta_{k}^{T}{(A}^{T}{A)\delta}_{k}) + {2\mathrm{\Delta}x}_{k}\left( {A\delta}_{k} \right)^{T}(AX - B) + \lambda\left| x_{k} + {\mathrm{\Delta}x}_{k} \right| + const\)

=
\(\left( {\mathrm{\Delta}x}_{k} \right)^{2}C_{k,k} + {2\mathrm{\Delta}x}_{k}A_{k}^{T}R + \lambda\left| x_{k} + {\mathrm{\Delta}x}_{k} \right| + const\)

\(f'\left( {\mathrm{\Delta}x}_{k} \right)\) =
2\({\mathrm{\Delta}x}_{k}C_{k,k} + 2A_{k}^{T}R + \mu_{k}\lambda\)

At optimum :

\(f'\left( {\mathrm{\Delta}x}_{k,opt} \right)\) =
2\({\mathrm{\Delta}x}_{k,opt}C_{k,k} + 2A_{k}^{T}R + \mu_{k}\lambda\) =
0

\({\mathrm{\Delta}x}_{k,opt}\) = -
\(\frac{A_{k}^{T}R}{C_{k,k}} - \frac{\mu_{k}\lambda}{{2C}_{k,k}}\)

\emph{where} \(A_{k}\) \emph{=} \({A\delta}_{k}\) \emph{= k-th column
of} \(A\) \emph{(column vector of size N)}

\(C_{k,k}\) \emph{=} \(\delta_{k}^{T}{(A}^{T}{A)\delta}_{k}\) \emph{=
(k,k) element of} \(C\) \emph{i.e.} \(A^{T}A\)

\emph{R =} \(AX - B\)

\(\mu_{k}\) = \(\left\lbrack \begin{matrix}
 + 1 & if & x_{k} + {\mathrm{\Delta}x}_{k} > 0 \\
 - 1 & if & x_{k} + {\mathrm{\Delta}x}_{k} < 0 \\
\lbrack - 1, + 1\rbrack & if & x_{k} + {\mathrm{\Delta}x}_{k} = 0
\end{matrix} \right.\ \) = \emph{slope of}
\(\left| x_{k} + {\mathrm{\Delta}x}_{k} \right|\)

\emph{case 1 :} \(x_{k} + {\mathrm{\Delta}x}_{k,opt} > 0\)\emph{,}
\(\mu_{k} = + 1\) \emph{then optimal} \({\mathrm{\Delta}x}_{k,opt}\)=-
\(\frac{A_{k}^{T}R}{C_{k,k}} - \frac{\lambda}{{2C}_{k,k}}\) \emph{must
fulfill condition :}

\({x_{k} + \mathrm{\Delta}x}_{k,opt}\) =
\(x_{k} - \frac{A_{k}^{T}R}{C_{k,k}} - \frac{\lambda}{2C_{k,k}}\)
\textgreater{} 0

\(x_{k} - \frac{A_{k}^{T}R}{C_{k,k}}\) \textgreater{}
+\(\frac{\lambda}{2C_{k,k}}\)

\emph{case 2 :}
\(x_{k} + {\mathrm{\Delta}x}_{k,opt} < 0\)\emph{,}\(\mu_{k} = - 1\)
\emph{then optimial} \({\mathrm{\Delta}x}_{k,opt}\)=-
\(\frac{A_{k}^{T}R}{C_{k,k}} + \frac{\lambda}{{2C}_{k,k}}\) \emph{must
fulfill condition :}

\({x_{k} + \mathrm{\Delta}x}_{k,opt}\) =
\(x_{k} - \frac{A_{k}^{T}R}{C_{k,k}} + \frac{\lambda}{2C_{k,k}}\)
\textless{} 0

\(x_{k} - \frac{A_{k}^{T}R}{C_{k,k}}\) \textless{}
\(- \frac{\lambda}{2C_{k,k}}\)

\emph{case 3 :} \(x_{k} + {\mathrm{\Delta}x}_{k,opt} = 0\)\emph{,}
\(\mu_{k} \in \lbrack - 1, + 1\rbrack\) \emph{then optimial}
\({\mathrm{\Delta}x}_{k,opt}\)=-\(x_{k}\) \emph{must fulfill condition
:}

\({x_{k} + \mathrm{\Delta}x}_{k,opt}\) =
\(x_{k} - \frac{A_{k}^{T}R}{C_{k,k}} - \frac{\mu_{k}\lambda}{2C_{k,k}}\)
= 0

\(x_{k} - \frac{A_{k}^{T}R}{C_{k,k}}\) =
\(- \frac{\mu_{k}\lambda}{2C_{k,k}}\) ∈ \(\pm \frac{\lambda}{2C_{k,k}}\)
\emph{(the kink)}\\
Given a positive semidefinite covariance \(C = A^{T}A\), quadratic
optimization must have one global optimum. For manually chosen penalty
weight \(\lambda\), \emph{X} is intermediate solution for previous
iteration, and \({\mathrm{\Delta}x}_{k,opt}\) is the optimal change in
one dimension for next iteration :

\emph{• if X and} \(x_{k}\) \emph{satisfy}
\(x_{k} - \frac{A_{k}^{T}R}{C_{k,k}} > + \frac{\lambda}{2C_{k,k}}\)

\emph{we can safely assume} \(x_{k} + {\mathrm{\Delta}x}_{k,opt} > 0\)
\emph{and get solution}
\({\mathrm{\Delta}x}_{k,opt} = - \frac{A_{k}^{T}R}{C_{k,k}} - \frac{\lambda}{2C_{k,k}}\)

\emph{i.e. taking simple least square solution of dimension k,
suppressed by} \(\frac{\lambda}{2C_{k,k}}\)

\emph{• if X and} \(x_{k}\) \emph{satisfy}
\(x_{k} - \frac{A_{k}^{T}R}{C_{k,k}} < - \frac{\lambda}{2C_{k,k}}\)

\emph{we can safetly assume} \(x_{k} + {\mathrm{\Delta}x}_{k,opt} < 0\)
\emph{and get solution}
\({\mathrm{\Delta}x}_{k,opt} = - \frac{A_{k}^{T}R}{C_{k,k}} + \frac{\lambda}{2C_{k,k}}\)

\emph{i.e. taking simple least square solution of dimension k, enhanced
by} \(\frac{\lambda}{2C_{k,k}}\)

\emph{• if X and} \(x_{k}\) \emph{lies in dead zone}
\(- \frac{\lambda}{2C_{k,k}} < x_{k} - \frac{A_{k}^{T}R}{C_{k,k}} < + \frac{\lambda}{2C_{k,k}}\)

\emph{both assumptions} \(x_{k} + {\mathrm{\Delta}x}_{k,opt} > 0\)
\emph{and} \(x_{k} + {\mathrm{\Delta}x}_{k,opt} < 0\) \emph{will be
wrong}

\emph{as this is a convex optimization with positive semidefinite C, the
only possible optimal is} \({\mathrm{\Delta}x}_{k,opt} = - x_{k}\)

In conclusion, optimum move can be found by soft-threshold operator :

\emph{if}
\(x_{k} - \frac{A_{k}^{T}R}{C_{k,k}} > + \frac{\lambda}{2C_{k,k}}\)
\emph{then}
\({\mathrm{\Delta}x}_{k,opt} = - \frac{A_{k}^{T}R}{C_{k,k}} - \frac{\lambda}{2C_{k,k}}\)

\emph{if}
\(x_{k} - \frac{A_{k}^{T}R}{C_{k,k}} < - \frac{\lambda}{2C_{k,k}}\)
\emph{then}
\({\mathrm{\Delta}x}_{k,opt} = - \frac{A_{k}^{T}R}{C_{k,k}} + \frac{\lambda}{2C_{k,k}}\)

\emph{if}
\(x_{k} - \frac{A_{k}^{T}R}{C_{k,k}} \in \pm \frac{\lambda}{2C_{k,k}}\)
\emph{then} \({\mathrm{\Delta}x}_{k,opt} = - x_{k}\) \emph{(dead zone,
making} \(x_{k}\) \emph{zero)}

Please note that :

1. dead zone is a region where previous value \emph{X} and \(x_{k}\)
lie, not where \({\mathrm{\Delta}x}_{k,opt}\) lies

2. dead zone makes the value of k-th parameter zero, that is
\({x'}_{k} = x_{k} + {\mathrm{\Delta}x}_{k,opt} = x_{k} - x_{k} = 0\)

3. dimension \emph{k} may lie in dead zone, while other dimensions may
not

4. dimension with \({x'}_{k} = 0\) implies feature is suppressed, LASSO
is effectively a feature selection

\emph{\ul{Comparison to SVM}}

• LASSO is a \emph{\textbf{feature selection}} algorithm, resulting in a
sparse set of features.

• SVM is a \emph{\textbf{data selection}} algorithm, resulting in sparse
set of data points.
